\section{Limitations and Future Work}

Monte-Carlo TopK speeds up selection; it does not replace indexing. In flat
search, it still needs the full score row. Its end-to-end gain depends on how
much time exact search spends in selection. A natural next step is to combine it
with product quantization, coarse clustering, or graph search.

The current implementation also pays a high cost for fallback. In SIFT1M, 18
underflow fallback rows make the full run slower than exact top-$K$. A
production implementation should use the same full-row selector as the exact
baseline, or retry only the rows with too few or too many candidates.

The current experiments report returned score, ID overlap, equal-score ties,
candidate tails, and tail latency. They still do not cover all $K$ values, batch
sizes, dimensions, data types, query shards, metrics, or indexed pipelines. The
results are for squared-L2 flat search. Inner product and cosine search should
be tested before making broader claims. A deployed system should choose the
threshold rank from online candidate statistics or from a calibration pass, not
only from a fixed table.

The implementation leaves several opportunities for optimization. It uses
separate kernels for sampled threshold selection, filtering/counting,
candidate compaction, and candidate top-$K$. A tighter implementation could fuse
adjacent stages, use prefix-sum compaction where it reduces atomic contention,
overlap sample scoring with full distance computation, and choose threshold
ranks per row or per batch. These changes could reduce overhead and fallback risk;
they do not change the main limitation that flat search still computes the full
score row.
